\section{Problem Formulation}

\subsection{System Dynamics}

Consider an $n$-degree-of-freedom (DOF) rigid robotic manipulator governed by the Euler--Lagrange equation

\begin{equation}
	M(q)\ddot q+C(q,\dot q)\dot q+G(q)=\tau+d,
	\label{eq:dynamics}
\end{equation}

where $q,\dot q,\ddot q\in\mathbb{R}^n$ denote the joint position, velocity, and acceleration vectors, respectively. The matrix $M(q)\in\mathbb{R}^{n\times n}$ denotes the inertia matrix, $C(q,\dot q)\in\mathbb{R}^{n\times n}$ is the Coriolis and centrifugal matrix, $G(q)\in\mathbb{R}^n$ is the gravity vector, $\tau\in\mathbb{R}^n$ is the control input, and $d\in\mathbb{R}^n$ represents the lumped external disturbance and unmodeled dynamics.

The following standard properties of Euler--Lagrange systems are adopted throughout this paper.

\begin{assumption}
	The unknown external disturbance $d(t)$ is piecewise continuous and uniformly bounded, i.e., there exists an unknown positive constant $d_m$ such that $\Vert{}d(t)\Vert{} \le d_m, \forall t \ge 0$. Furthermore, if the disturbance depends on the system states $d(x, t)$, it is assumed to be Lipschitz continuous with respect to $x$ over a compact domain $\Omega_x$, ensuring the existence of a bounded RBF approximation error.
\end{assumption}

\begin{property}
	The inertia matrix $M(q)$ is symmetric and uniformly positive definite. Specifically, there exist two positive constants $m_1$ and $m_2$ satisfying
	\begin{equation}
		m_1I\le M(q)\le m_2I,
	\end{equation}
	for all $q\in\mathbb{R}^n$.
\end{property}

\begin{property}
	The matrix $\dot M(q)-2C(q,\dot q)$ is skew-symmetric, i.e.,
	\begin{equation}
		x^{\mathrm T}\!\left(\dot M(q)-2C(q,\dot q)\right)x=0,
	\end{equation}
	for any $x\in\mathbb{R}^n$.
\end{property}

In practical robotic systems, the exact dynamic model is unavailable due to parameter uncertainties, payload variations, friction, and unmodeled dynamics. Accordingly, the manipulator dynamics are decomposed into a nominal model and an uncertainty term as

\begin{equation}
	M(q)=M_0(q)+\Delta M(q),
\end{equation}

\begin{equation}
	C(q,\dot q)=C_0(q,\dot q)+\Delta C(q,\dot q),
\end{equation}

\begin{equation}
	G(q)=G_0(q)+\Delta G(q),
\end{equation}

where $M_0(q)$, $C_0(q,\dot q)$, and $G_0(q)$ denote the nominal inertia, Coriolis/centrifugal, and gravity terms, respectively.

Substituting the above decompositions into \eqref{eq:dynamics} yields

\begin{equation}
	\left(M_0+\Delta M\right)\ddot q
	+\left(C_0+\Delta C\right)\dot q
	+\left(G_0+\Delta G\right)
	=
	\tau+d.
\end{equation}

Rearranging the above equation gives

\begin{equation}
	M_0(q)\ddot q
	+
	C_0(q,\dot q)\dot q
	+
	G_0(q)
	=
	\tau
	+
	f(q,\dot q,\ddot q),
	\label{eq:nominal}
\end{equation}

where

\begin{equation}
	f(q,\dot q,\ddot q)
	=
	d
	-
	\Delta M(q)\ddot q
	-
	\Delta C(q,\dot q)\dot q
	-
	\Delta G(q).
	\label{eq:lumped}
\end{equation}

The nonlinear function $f(\cdot)$ represents the lumped uncertainty consisting of modeling errors, parameter variations, external disturbances, and unmodeled dynamics.

%%%%%%%%%%%%%%%%%%%%%%%%%%%%%%%%%%%%%%%%%%%%%%%%%%%%%%%%%%%%%%

\subsection{Tracking Control Problem}

Assume that the desired joint trajectory satisfies

\begin{equation}
	q_d,\;
	\dot q_d,\;
	\ddot q_d
	\in
	\mathcal{L}_{\infty}.
\end{equation}

Define the tracking error and its derivative as

\begin{equation}
	e=q-q_d,
	\qquad
	\dot e=\dot q-\dot q_d.
	\label{eq:error}
\end{equation}

Differentiating \eqref{eq:error} yields

\begin{equation}
	\ddot e=\ddot q-\ddot q_d.
	\label{eq:error_second}
\end{equation}

The desired tracking performance is specified by the second-order error dynamics

\begin{equation}
	\ddot e
	+
	K_v\dot e
	+
	K_pe
	=
	0,
	\label{eq:error_dynamics}
\end{equation}

where $K_p,K_v\in\mathbb{R}^{n\times n}$ are symmetric positive-definite gain matrices.

\begin{assumption}
	The external disturbance is bounded, i.e.,
	\begin{equation}
		\|d(t)\|\le d_{\max},
	\end{equation}
	where $d_{\max}$ is an unknown positive constant.
\end{assumption}

Substituting \eqref{eq:error_second} into \eqref{eq:error_dynamics} yields the desired joint acceleration

\begin{equation}
	\ddot q
	=
	\ddot q_d
	-
	K_v\dot e
	-
	K_pe.
	\label{eq:desired_acc}
\end{equation}

The control objective is to design an adaptive neural controller such that all closed-loop signals remain bounded and the tracking error converges to a sufficiently small neighborhood of the origin despite the presence of lumped uncertainties.